\section{Target--coupled CRT faces as character projectors}
\label{sec:tpc38-character-projectors}

We first work with any fixed number \(r\ge1\) of distinct odd primes.
Put
\begin{equation}
 M=\prod_{i=1}^rq_i,\qquad
 n_i=q_i-1,\qquad p_i=n_i^{-1},\qquad
 \Phi=\varphi(M)=\prod_{i=1}^rn_i.
 \label{eq:tpc38-general-M}
\end{equation}
Every Dirichlet character modulo \(M\) is viewed as a tensor product of
local multiplicative characters and is extended by zero away from the
unit group.  Throughout this section, \(N\in(\Z/M\Z)^\times\).

\subsection{Coordinate faces}

For \(F\in\Z/M\Z\), define locally
\begin{align}
 d_{i,N}(F)&=\one_{F=0\ (\bmod q_i)},
 \label{eq:tpc38-local-d}\\
 e_{i,N}(F)&=\one_{N+F\ne0\ (\bmod q_i)},
 \label{eq:tpc38-local-e}\\
 m_{i,N}(F)&=p_i e_{i,N}(F),
 \qquad
 c_{i,N}(F)=d_{i,N}(F)-m_{i,N}(F).
 \label{eq:tpc38-local-m-c}
\end{align}
The exclusion in \eqref{eq:tpc38-local-e} is forced by zero extension of
multiplicative characters.  For \(S\subseteq[r]\), put
\begin{equation}
 K_{S,N}(F)
 =\prod_{i\in S}c_{i,N}(F)
  \prod_{i\notin S}m_{i,N}(F).
 \label{eq:tpc38-coordinate-face}
\end{equation}

Let \(\widehat G_i\) be the characters of
\(G_i=\F_{q_i}^\times\), including the principal character
\(\mathbf1_i\).  Define
\begin{equation}
 \mathfrak X_S
 =\left\{\chi=\bigotimes_{i=1}^r\chi_i:
 \chi_i\ne\mathbf1_i\ \Longleftrightarrow\ i\in S\right\}.
 \label{eq:tpc38-character-family-S}
\end{equation}

\begin{theorem}[Exact target--coupled character projector]
\label{thm:tpc38-character-projector}
For every \(S\subseteq[r]\), every unit \(N\bmod M\), and every
\(F\bmod M\),
\begin{equation}
 \boxed{
 K_{S,N}(F)
 =\frac1\Phi\sum_{\chi\in\mathfrak X_S}
 \chi(N+F)\overline{\chi(N)}.}
 \label{eq:tpc38-character-projector}
\end{equation}
If \(N+F=su\), then each summand factors as
\begin{equation}
 \chi(N+F)\overline{\chi(N)}
 =\chi(s)\chi(u)\overline{\chi(N)}.
 \label{eq:tpc38-character-product-factor}
\end{equation}
Both identities remain valid when \(N+F\), \(s\), or \(u\) is a nonunit.
\end{theorem}

\begin{proof}
In one local coordinate, multiplicative orthogonality gives
\begin{equation}
 d_{i,N}(F)
 =\frac1{n_i}\sum_{\chi_i\in\widehat G_i}
 \chi_i(N+F)\overline{\chi_i(N)}.
 \label{eq:tpc38-local-all-character}
\end{equation}
If \(N+F=0\), both sides vanish.  The principal term on the right is
\(m_{i,N}(F)\), because \(N\) is a unit.  Subtracting it gives
\begin{equation}
 c_{i,N}(F)
 =\frac1{n_i}\sum_{\substack{\chi_i\in\widehat G_i\\
 \chi_i\ne\mathbf1_i}}
 \chi_i(N+F)\overline{\chi_i(N)}.
 \label{eq:tpc38-local-nonprincipal}
\end{equation}
Tensor \eqref{eq:tpc38-local-all-character} and
\eqref{eq:tpc38-local-nonprincipal} through the CRT to prove
\eqref{eq:tpc38-character-projector}.  Complete multiplicativity gives
\eqref{eq:tpc38-character-product-factor}.  If one factor is a nonunit,
the relevant character values vanish on both sides.
\end{proof}

For \(r=3\), define
\begin{equation}
 \mathfrak X_{\rm full}=\mathfrak X_{[3]},
 \qquad
 \mathfrak X_{\rm prop}
 =\bigcup_{S\subsetneq[3]}\mathfrak X_S,
 \label{eq:tpc38-full-proper-families}
\end{equation}
and
\begin{equation}
 C_N=K_{[3],N},\qquad
 P_N=\sum_{S\subsetneq[3]}K_{S,N}.
 \label{eq:tpc38-full-proper-faces}
\end{equation}
With
\begin{equation}
 A=\prod_{i=1}^3(1-p_i),
 \label{eq:tpc38-A-definition}
\end{equation}
their family sizes are
\begin{equation}
 |\mathfrak X_{\rm full}|=A\Phi,
 \qquad
 |\mathfrak X_{\rm prop}|=(1-A)\Phi.
 \label{eq:tpc38-family-sizes}
\end{equation}

\begin{corollary}[Full/proper recombination]
\label{cor:tpc38-face-recombination}
For every unit target \(N\) modulo \(M\),
\begin{equation}
 \boxed{C_N(F)+P_N(F)=\one_{M\mid F}.}
 \label{eq:tpc38-face-recombination}
\end{equation}
In particular,
\begin{equation}
 C_N(0)=A,\qquad P_N(0)=1-A.
 \label{eq:tpc38-face-values-zero}
\end{equation}
If \(M\nmid F\), then
\begin{equation}
 P_N(F)=-C_N(F).
 \label{eq:tpc38-offjoint-cancellation}
\end{equation}
\end{corollary}

\begin{proof}
The two families in \eqref{eq:tpc38-full-proper-families} partition all
characters modulo \(M\).  Apply character orthogonality to their sum.  At
\(F=0\), all character phases equal one, so the two values are their
family densities.  The final identity follows from the first.
\end{proof}

\subsection{Fixed target Parseval and shifted target Gram}

For a fixed unit \(N\), put
\begin{equation}
 e_{\chi,N}(F)
 =\chi(N+F)\overline{\chi(N)}.
 \label{eq:tpc38-echiN}
\end{equation}

\begin{proposition}[Fixed target Parseval]
\label{prop:tpc38-fixed-target-Parseval}
For characters \(\chi,\psi\bmod M\),
\begin{equation}
 \sum_{F\bmod M}e_{\chi,N}(F)
 \overline{e_{\psi,N}(F)}
 =\Phi\one_{\chi=\psi}.
 \label{eq:tpc38-fixed-target-orthogonality}
\end{equation}
Consequently,
\begin{equation}
 \boxed{
 \sum_{F\bmod M}|C_N(F)|^2=A,\qquad
 \sum_{F\bmod M}|P_N(F)|^2=1-A.}
 \label{eq:tpc38-fixed-target-face-norms}
\end{equation}
\end{proposition}

\begin{proof}
The translation \(F\mapsto N+F\) runs through every residue modulo \(M\).
Character orthogonality on the unit group proves
\eqref{eq:tpc38-fixed-target-orthogonality}.  Insert the disjoint character
families and their sizes from \eqref{eq:tpc38-family-sizes}.
\end{proof}

When the target moves, the Gram is no longer diagonal.  For units \(N,N'\)
let \(\Delta=N'-N\) and define
\begin{equation}
 \Gamma_{\chi,\psi}(N,N')
 =\sum_{F\bmod M}e_{\chi,N}(F)
 \overline{e_{\psi,N'}(F)}.
 \label{eq:tpc38-shifted-Gram-definition}
\end{equation}

\begin{theorem}[Shifted Jacobi Gram]
\label{thm:tpc38-shifted-Jacobi-Gram}
With local components denoted by subscripts,
\begin{equation}
 \boxed{
 \Gamma_{\chi,\psi}(N,N')
 =\overline{\chi(N)}\psi(N')
 \prod_{i=1}^r
 S_i(\chi_i,\psi_i;\Delta_i),}
 \label{eq:tpc38-shifted-Jacobi-Gram}
\end{equation}
where
\begin{equation}
 S_i(\chi_i,\psi_i;\Delta_i)
 =\sum_{x\bmod q_i}\chi_i(x)
 \overline{\psi_i(x+\Delta_i)}.
 \label{eq:tpc38-local-shifted-sum}
\end{equation}
If \(\Delta_i=0\), then
\begin{equation}
 S_i=n_i\one_{\chi_i=\psi_i}.
 \label{eq:tpc38-local-shifted-zero}
\end{equation}
If \(\Delta_i\ne0\), then
\begin{equation}
 \boxed{
 S_i
 =\chi_i(-\Delta_i)\overline{\psi_i(\Delta_i)}
 J(\chi_i,\overline{\psi_i}),}
 \label{eq:tpc38-local-shifted-Jacobi}
\end{equation}
where \(J(\alpha,\beta)=\sum_x\alpha(x)\beta(1-x)\).
\end{theorem}

\begin{proof}
Set \(x=N+F\).  Then \(N'+F=x+\Delta\), and the CRT factors the sum in
\eqref{eq:tpc38-shifted-Gram-definition} into
\eqref{eq:tpc38-local-shifted-sum}.  The zero--shift case is ordinary
character orthogonality.  If \(\Delta_i\ne0\), substitute
\(x=-\Delta_i y\).  Then \(x+\Delta_i=\Delta_i(1-y)\), which gives
\eqref{eq:tpc38-local-shifted-Jacobi}.
\end{proof}

\begin{remark}[What has and has not been diagonalized]
\label{rem:tpc38-fixed-versus-moving-target}
The moving affine determinant has been converted into an exact ratio
phase.  The price is the shifted Jacobi Gram
\eqref{eq:tpc38-shifted-Jacobi-Gram}; fixed--target Parseval does not
persist when target rows are synthesized coherently.  This distinction is
the finite spectral content of the target--coupling problem.
\end{remark}
